% Section 7 -- Case Study
\section{Case Study: AS12479 Route Leak}
\label{sec:casestudy}

We trace one event end-to-end to make the paper's
contributions concrete. The AS12479 route leak of 2024-01-03
exercises all four questions (Q0--Q3), is backed by a primary-source
URL, and has a complete cached run: routing detections, a
PathProb path-probability snapshot, a CAIDA AS-relationship
snapshot, and Cloudflare Radar traffic.

    10|\medskip
\noindent\textbf{What happened.} Between 14:45\,UTC and 18:00\,UTC
on 2024-01-03 the operational AS12479 announced prefixes that
PathProb's path-probability model scored as highly improbable given
the AS-relationship graph. The leak detector scored 13{,}716 candidate
announcements against 1{,}886{,}291 announcements observed in the
window, of which 14{,}776 survived the AS-path pre-filter, and
reported the 5{,}000 highest-scoring as leaks, capped by a
truncation limit. The leak was scored against a 536{,}283-entry
AS-relationship snapshot at a path-probability threshold of 0.4.
The same window also produced a candidate outage signal with a
low raw confidence of 0.3.

\medskip
\noindent\textbf{Stage 1 -- Coarse screen.} The week-over-week
baseline comparison produced a quick score below the 0.20
admission threshold, because the Cloudflare Radar series for
AS12479 had no usable periodic baseline in this window:
the baseline method fell back to a non-periodic estimate and
the periodicity confidence was 0.0. The traffic detector still
reported a peak-to-trough excursion of 81.7\%, but with zero
discrete anomalies and an empty historical average, so the
screen did not admit the candidate on traffic grounds alone.
The routing evidence was what carried the event.

    30|\medskip
\noindent\textbf{Phase 1 -- Deterministic arbitration.}
Each detector produced a structured argument:

\begin{itemize}
  \item \textbf{Leak}: fires strongly. 13{,}716 candidate
    announcements scored above the 0.4 path-probability threshold;
    the leak detector returned a saturated 5{,}000-record result.
    Because the result was truncated, the count in the report
    is a floor rather than an estimate.
  \item \textbf{Traffic}: fires weakly. The 81.7\% peak-to-trough
    excursion is a real signal, but with no periodic baseline and
    no discrete anomalies the detector cannot place it in time,
    so it cannot be aligned against the routing event and the
    traffic channel is suspended at the Phase-1 gate.
  \item \textbf{Outage}: fires weakly. A withdrawal signal was
    produced, but with a raw confidence of 0.3 it falls below the
    evidential gate.
  \item \textbf{Origin-hijack}: silent. No origin hijack was
    detected in this window; the detector is suspended, not
    silenced, and its reason is recorded in the cache.
\end{itemize}
The arbitration therefore returns \texttt{silent\_alert} for the
Phase-1 correlation verdict: a routing signal is present
while the traffic signal is absent, so the Phase-1 arbitration
suspends the traffic channel without reaching the temporal
correlation test. On the evidence available in this window the
routing anomaly cannot be shown to have reached traffic. This is
the correct Phase-1 outcome, and it is exactly the verdict that
tells an operator to escalate rather than close. A routing-only
pipeline would have stopped here with a bare leak count; a
traffic-only pipeline would never have seen the event at all.

    65|\medskip
\noindent\textbf{Phase 2 -- Pearl-PNS causal attribution.}
The Bayesian synthesis committed to $R \to T$ with a posterior of
0.975 on $R \to T$, 0.011 on $T \to R$, and 0.014 on independent,
against an independently verified ground truth of the
$R \to T$ direction. This is the case that motivates the
abstention gate. Every individual classical channel failed on this
event: the CSS estimator abstained because the bootstrap CI could
not exclude zero over a series with too few non-zero buckets,
Granger abstained on the bursty update stream, and Pearson
correlation abstained as well. The single transfer-entropy
estimate pointed the wrong way. Only the synthesis, which discounts
a channel that disagrees with its Pearl-anchored prior instead of
letting it vote, produced the correct direction. The event also
illustrates the truncation caveat: the leak count that reaches the
report is a lower bound, so the evidence score attached to the
leak detector is a floor, not a point estimate.

    85|\medskip
\noindent\textbf{Why not a single detector?} An
outage-detector-only system sees a low-confidence withdrawal
signal and stops. A leak-only system reports the leak but cannot
say whether it had traffic impact, which is the operator's
actual question. Only BGP Tracer answers Q0--Q3 together: the
event is real, the leak is the cause class, and the causal
direction is committed with a stated posterior.

    95|\medskip
\noindent\textbf{Operator loop and latency.} Triage is
automatic; the cached Phase-1 run for this event completed in
273\,s end-to-end, dominated by the PathProb and
AS-relationship snapshot loads rather than by the causal engine.
The Phase-2 attribution adds a median of 8.2\,s when invoked.

\begin{figure}[t]
\centering
\includegraphics[width=0.92\linewidth]{figs/fig_case_as12479.pdf}
\caption{AS12479 route leak, 2024-01-03.
Top: Cloudflare Radar five-minute traffic for the week; the shaded
band is the analysis window, in which the series falls from a peak of
0.96 to a trough of 0.17.
Middle: yield of the PathProb leak detector, from 1.9\,M observed
announcements down to the 5{,}000 reported leaks, which is a truncated
count and therefore a floor.
Bottom: the Bayesian posterior over causal directions. The dominant
$R \to T$ hypothesis at 0.975 clears the commit floor $\tau=0.34$,
even though every individual classical channel abstained.}
\label{fig:case}
\end{figure}

    115|The other two anomaly classes are illustrated with condensed
timelines in Appendix~\ref{app:case-geo}: a Cloudflare
data-centre outage of 2020-04-15 and a Swiss operator route
leak of 2019-06-05. The detector-output geometry for all
curated 2025 events is in Appendix~\ref{app:case-geo}.
